\chapter{Introduction}
\label{chap:introduction}

% Background on urban traffic congestion at the Thapathali intersection in Kathmandu
\section{Background}
\label{sec:background}

Kathmandu's congestion is most visible at intersections, where several
busy approaches compete for limited green time. Thapathali is a suitable
case because traffic enters from Maitighar, Kupondole, Tripureshwor, and
Maternity, and the dominant movement changes during the day. In practice,
signal operation still depends on fixed timing plans, camera observation,
and manual intervention by traffic police.

Those methods are understandable and easy to operate, but they react
slowly when a queue forms suddenly or when one peak-hour approach becomes
heavier than the others. For that reason, this report studies
intersection-level signal control at Thapathali instead of attempting to
model the full Kathmandu network.

% Rationale for applying deep reinforcement learning to traffic signal optimization
\section{Motivation}
\label{sec:motivation}

The motivation for using \gls{drl} comes from a practical limitation of
fixed timing. A fixed plan gives each phase a pre-decided share of time,
even when the queue distribution has changed. At Thapathali, this can
mean serving a lightly loaded direction while another approach continues
to build a queue.

A learning-based controller is useful to test because it does not need a
hand-written rule for every queue pattern. It observes the current state,
applies a phase decision in simulation, and updates its policy from the
delay and queue behaviour that follows.

This work uses \gls{sumo} because real signal experiments at a busy
Kathmandu junction would be unsafe and outside the scope of a minor
project. Simulation allows the same Day 5 demand to be tested under a
fixed 60-second plan and under the trained DQN policy. The purpose is not
to claim immediate field deployment; it is to check whether the adaptive
controller gives a measurable improvement in a controlled version of the
intersection.

% Definition of the core problem and the DQN-based hypothesis to address it
\section{Problem Statement}
\label{sec:problem_statement}

The main problem addressed in this project is the mismatch between fixed signal timing and changing traffic demand at Thapathali. A controller that does not observe the current queue condition may hold a phase longer than necessary or switch too late when another approach is already congested. This increases waiting time, queue length, and the chance of spillback into nearby links.

The project hypothesis is that a \gls{dqn}-based controller can reduce
average waiting time and queue length compared with fixed-time control
under the same simulated demand. The expected advantage comes from using
PCU-weighted lane loads, stopline waiting time, and the current phase
rather than relying only on pre-set splits.

The expected benefit is not unlimited. If all approaches are saturated at
the same time, the controller can only distribute scarce green time more
carefully; it cannot create additional road capacity.

% Primary and measurable goals of the project
\section{Objectives}
\label{sec:objectives}

The objectives were defined so that the controller could be evaluated against a baseline rather than described only at a conceptual level:

\begin{itemize}
    \item To prepare a \gls{sumo} simulation of the Thapathali intersection using available road geometry, OD data, vehicle counts, and PCU-based vehicle classes.
    \item To develop a \gls{dqn}-based \gls{drl} agent with a defined state vector, action space, reward function, and training process for adaptive signal control.
    \item To compare the trained agent with a fixed-time signal baseline using average waiting time, queue length, and throughput efficiency.
\end{itemize}

% Boundaries of the project and known constraints on the implementation
\section{Scope and Applications}
\label{sec:scope_limitations}

This project is limited to a single-intersection simulation of
Thapathali. It does not control the wider Kathmandu network, and it does
not include hardware deployment at a real signal controller. Within that
boundary, the work covers network preparation, traffic-demand generation,
DQN training, and comparison with a fixed-time baseline. The results are
therefore best read as a simulation study and a starting point for later
traffic-engineering experiments.

\subsection{Scope}
\label{subsec:scope}

The main scope items are:

\begin{itemize}
    \item Configuring the Thapathali intersection map in \gls{sumo}, including lane layout, signal phases, and vehicle routes.
    \item Designing the DQN controller around the selected state representation, STAY/NEXT action space, and Max-Pressure based reward.
    \item Evaluating the trained policy against a fixed 60-second signal plan using average waiting time, queue length, and throughput efficiency.
\end{itemize}
